% limitations_conclusion.tex  -- written by Shivansh
\subsection{Limitations}
The FPGA numbers come from an open-source flow: post-route on Lattice ECP5,
and synthesis-level estimates for Xilinx 7-series. Results from the
vendor's own tools for a Xilinx device may differ by tens of percent. The sequential baseline was not pipelined, so its 75.7~MHz
understates what a tuned single-MAC design could reach; the cycle-based
speed-ups (5.4$\times$ compute, 9.4$\times$ end-to-end) do not depend on
this. Finally, the OS design reloads both operands for every product,
which limits utilisation to 1/3, as discussed above.

% -----------------------------------------------------------------------------
\section{Conclusion}
We designed, verified and evaluated a parameterised output-stationary
systolic-array accelerator. The Phase~I prototype showed that the array
computes correctly with a $3N$-cycle latency. Phase~II showed that the
array alone does not determine performance. Data transfer initially
consumed 80\,\% of the time, and only row-wide transfers, double buffering
and a result buffer turned the 5.4$\times$ compute advantage into a real
end-to-end gain. Two pipeline registers then doubled the clock rate. The
final $4\times4$ design completes a product every 14 cycles at 116~MHz
(121~ns): 12.4$\times$ faster than a single-MAC baseline, verified by more
than 85\,000 checks and predicted exactly by a simple analytical model.
The weight-stationary comparison quantified why large accelerators such as
the TPU keep weights in place: utilisation rises from 1/3 towards 1 when
operands are reused across many rows. Natural extensions are a
weight-stationary or hybrid controller with on-chip activation buffers, a
DMA-style memory interface, deeper DSP-internal pipelining and running the
design on an FPGA board.
